\documentclass[11pt,a4paper]{article}
\usepackage[utf8]{inputenc}
\usepackage[T1]{fontenc}
\usepackage[spanish,es-tabla]{babel}
\usepackage{lmodern}
\usepackage[margin=2.2cm]{geometry}
\usepackage{booktabs}
\usepackage{amsmath}
\usepackage{enumitem}
\usepackage{parskip}
\setlist{nosep,leftmargin=*}
\renewcommand{\arraystretch}{1.1}

\title{\vspace{-1.5cm}Distancia al óptimo de la mejora por grupos}
\author{Capstone --- buses eléctricos}
\date{Reunión con el profesor, 16 de noviembre de 2026}

\begin{document}
\maketitle
\vspace{-0.6cm}

\textbf{Resultado.} En los 5 casos chicos donde Gurobi prueba el óptimo, el método que
usamos en la red completa (cuartetos al azar) queda en promedio a \textbf{0,13\,\%} del
costo óptimo, y a 0,31\,\% en el peor caso. Usa siempre el mismo número de buses que el
óptimo. El plan inicial, sin mejora, quedaba a 1,39\,\%.

\section{Qué queríamos saber}
Nuestro método parte de un plan inicial y lo mejora reasignando viajes entre grupos de
buses. Es una heurística: no garantiza el óptimo. En la red completa (417 rutas, unos
7.350 buses en día laboral) el modelo exacto no se puede resolver, así que no sabemos
cuánto nos falta. Para estimarlo, armamos casos más chicos donde el óptimo sí se puede
calcular y comparamos.

\section{Cómo lo medimos}
\begin{enumerate}
  \item \textbf{Casos.} Tomamos subconjuntos de la instancia real, un sector a la vez:
    5 casos \emph{chicos} de 2 a 4 rutas (84 a 157 corridas) y 6 \emph{muy chicos} de
    una ruta (24 a 26 corridas).
  \item \textbf{Óptimo.} Un solo MILP en Gurobi decide qué bus hace cada corrida, dónde
    termina la jornada y cuándo recarga, con el mismo costo (USD 250 por bus, USD 0,5
    por km y USD 5 por sesión de carga) y el mismo catálogo de recargas que la
    heurística. Con tolerancia 0 y límite de 30 minutos, probó el óptimo en los 5 casos
    chicos, en 1 a 8 minutos cada uno.
  \item \textbf{Métodos.} Corrimos el plan inicial, pares, pares + tríos, la mezcla
    creciente (tipo VNS) y grupos de 3, 4, 5 y 8 buses, elegidos al azar (3 semillas) o
    por parecido. Cada método tuvo 60 y 120 segundos; ambos tiempos dan casi lo mismo.
  \item \textbf{Brecha.} Para cada corrida,
    \[
      \text{brecha} = \frac{\text{costo del método} - \text{costo óptimo}}{\text{costo óptimo}}.
    \]
    Promediamos primero las semillas de cada caso y después los 5 casos.
\end{enumerate}

\section{Resultados}
La Tabla~\ref{tab:casos} muestra cada caso chico y la Tabla~\ref{tab:metodos} resume
la brecha de cada método con 120 segundos.

\begin{table}[h]
\centering
\small
\caption{Casos chicos: óptimo probado y brecha de tres métodos.}
\label{tab:casos}
\begin{tabular}{lcccrrrrr}
\toprule
 & & & & \multicolumn{2}{c}{Óptimo} & \multicolumn{3}{c}{Brecha al óptimo (\%)} \\
\cmidrule(lr){5-6}\cmidrule(lr){7-9}
Caso & Sector & Rutas & Corridas & Buses & USD & Plan inicial & Pares + tríos & Cuartetos \\
\midrule
chico\_02 & 2 & 2 & 84  & 29 & 8.651,57  & 2,90 & 2,70 & 0,03 \\
chico\_03 & 3 & 3 & 116 & 24 & 7.430,86  & 3,52 & 0,25 & 0,19 \\
chico\_04 & 4 & 3 & 152 & 42 & 13.041,50 & 0,27 & 0,00 & 0,13 \\
chico\_05 & 5 & 4 & 157 & 35 & 11.225,71 & 0,12 & 0,06 & 0,12 \\
chico\_06 & 4 & 4 & 141 & 28 & 8.425,65  & 0,16 & 0,06 & 0,16 \\
\midrule
Promedio  &   &   &     &    &           & 1,39 & 0,61 & 0,13 \\
\bottomrule
\end{tabular}

\vspace{2pt}
{\footnotesize Cuartetos: grupos de 4 buses al azar, promedio de 3 semillas.}
\end{table}

\begin{table}[h]
\centering
\small
\caption{Brecha al óptimo por método en los 5 casos chicos (120 segundos).}
\label{tab:metodos}
\begin{tabular}{lccc}
\toprule
Método & Brecha media (\%) & Peor caso (\%) & Buses de más (promedio) \\
\midrule
Plan inicial, sin mejora   & 1,39 & 3,52 & 0,4 \\
Pares                      & 1,24 & 3,39 & 0,4 \\
Pares + tríos              & 0,61 & 2,70 & 0,2 \\
Mezcla creciente           & 0,56 & 2,70 & 0,2 \\
Tríos al azar              & 0,15 & 0,38 & 0 \\
\textbf{Cuartetos al azar} & \textbf{0,13} & \textbf{0,31} & \textbf{0} \\
Quintetos al azar          & 0,14 & 0,32 & 0 \\
Octetos al azar            & 0,14 & 0,38 & 0 \\
\bottomrule
\end{tabular}
\end{table}

Los métodos que parten con pares quedan atrapados en chico\_02: se quedan en 30 buses
cuando el óptimo usa 29, y los tríos no logran salir de ese plan. Los grupos de 3 o más
buses llegan a 29 desde el comienzo. Por eso la brecha baja de 0,61\,\% a 0,13\,\% al
pasar de pares + tríos a cuartetos.

En la red completa, el mismo método con 40 minutos baja el costo frente al plan inicial
en 2,25\,\% en día laboral, 4,80\,\% en sábado y 8,00\,\% en domingo.

\section{Qué podemos afirmar y qué no}
\begin{itemize}
  \item \textbf{Sí:} en casos de hasta 157 corridas, la mejora por grupos cierra cerca
    del 90\,\% de la distancia entre el plan inicial y el óptimo, y nunca usa buses de
    más.
  \item \textbf{No:} la brecha de la red completa. Los casos chicos tienen como máximo
    157 corridas; la red completa es mucho más grande y allí no tenemos cota inferior.
  \item \textbf{La ventaja de los cuartetos depende de un caso.} En chico\_04, 05 y 06,
    pares + tríos queda más cerca del óptimo (0,00 a 0,06\,\% contra 0,12 a 0,16\,\%), y
    los cuartetos no mejoran el plan inicial en chico\_05 ni en chico\_06.
  \item \textbf{Casos muy chicos (una ruta).} Pares y pares + tríos llegan al óptimo en
    los 6. En la ruta J22, los cuartetos quedan a 4,2\,\%. Los grupos de 5 u 8 no se
    pueden formar con los 4 buses del plan inicial y se quedan en el plan inicial (25,4\,\%).
  \item \textbf{Casi no hay recargas.} El óptimo no recarga en ninguno de los 5 casos
    chicos. La brecha mide sobre todo cómo se encadenan los viajes, no las decisiones
    de recarga.
  \item \textbf{Casos medianos (200 a 400 corridas): no se pudieron resolver.} Con 200
    corridas, el modelo exacto usó unos 13~GB de memoria en un computador de 8~GB.
\end{itemize}

\section{Preguntas para la reunión}
\begin{enumerate}
  \item ¿Basta con la brecha en casos chicos, o conviene buscar una cota inferior
    para la red completa (por ejemplo, la relajación lineal por sector)?
  \item ¿Vale la pena resolver los casos medianos en un equipo con más memoria?
  \item ¿Nos quedamos con cuartetos al azar, o combinamos pares + tríos con grupos más
    grandes, ya que cada uno gana en casos distintos?
\end{enumerate}

\end{document}
